\section{Conclusion}

Our research sheds light on the emerging ecosystem of CNAME-based tracking, a tracking scheme that takes advantage of a DNS-based cloaking technique to evade tracking countermeasures.
Using HTTP Archive data and a novel method, we performed a longitudinal analysis of the CNAME-based tracking ecosystem using crawl data of 5.6M web pages.
Our findings show that unlike other trackers with similar scale, CNAME-based trackers are becoming increasingly popular, and are mostly used to supplement ``typical'' third-party tracking services.
We evaluated the privacy and security threats that are caused by including CNAME trackers in a same-site context.
Through manual analysis we found that sensitive information such as email addresses and authentication cookies leak to CNAME trackers on sites where users can create accounts.
Furthermore, we performed an automated analysis of cookie leaks to CNAME trackers and found that cookies set by other parties leak to CNAME trackers on 95\% of the websites that we studied.
Finally we identified two major web security vulnerabilities that CNAME trackers caused.
We disclosed the vulnerabilities to the respective parties and have worked with them to mitigate the issues.
We hope that our research helps with addressing the security and privacy issues that we highlighted, and inform development of countermeasures and policy making with regard to online privacy and tracking. 
